\section{Related Work}

\subsection{Weight-Space Learning}

The treatment of neural network weights as a learnable data modality has emerged as an active research direction. The 2026 Weight Space Learning Survey categorizes this field into understanding, representation, and generation paradigms~\cite{han2026survey}. Our work falls under \emph{understanding}: decoding behavioral properties from weight inspection.

\paragraph{Permutation Equivariance.} A fundamental challenge in processing weights is that hidden neurons have no canonical ordering. Zhou et al.\ introduced Neural Functionals (NF-Layers) that respect this symmetry through parameter-sharing schemes~\cite{zhou2023nf}. Their characterization theorem proves that NF-Layers capture all permutation-equivariant linear maps.

\paragraph{Scalable Representations.} SANE addresses scalability by tokenizing weight subsets sequentially~\cite{schurholt2024sane}. ProbeGen uses learned deep linear probes with shared generators~\cite{kahana2024probegen}. Set-based Neural Network Encoding handles mixed-architecture model zoos through logit invariance~\cite{andreis2023setbased}.

\subsection{Property Prediction from Weights}

Prior work on property prediction has focused on scalar targets---primarily accuracy and generalization metrics.

\paragraph{Weight Statistics.} Unterthiner et al.\ demonstrated that handcrafted weight statistics achieve $R^2 \approx 0.97$ for accuracy prediction on the Small CNN Zoo~\cite{unterthiner2020predicting}. This established that scalar accuracy is highly predictable from simple features. Our work tests whether this extends to structured behavioral predictions.

\paragraph{Learned Representations.} SANE achieves $R^2 > 0.9$ for accuracy via learned embeddings~\cite{schurholt2024sane}. Herrmann et al.\ adapt Deep Weight Space layers for RNNs using ``functionalist interrogation''~\cite{herrmann2024rnn}.

\paragraph{The Gap.} These methods predict \emph{scalar} properties. Class-wise accuracy profiles represent structured behavioral predictions with cross-class dependencies. Whether behavioral information extends to structured profiles remains untested.

\subsection{Behavioral Understanding}

\paragraph{Behavioral Loss.} Meynent et al.\ propose behavioral loss for weight-space autoencoders, showing that structural reconstruction does not guarantee behavioral fidelity~\cite{meynent2025behavior}. They report 16--20\% accuracy drops despite low reconstruction error, motivating explicit behavioral supervision.
